\emph{(i).} By \claimref{std-algebra.complex-numbers}(ii), $J = e_{1}e_{2}$ has $J^{2} = -1$, and $\Z[J]$ is a commutative ring, the integer combinations of $1$ and $J$. Then $(J - 1)^{2} = J^{2} - 2J + 1 = -2J$ and $J(J - 1)^{2} = -2J^{2} = 2$. The reverse sends $J$ to $-J$, so $z\tilde{z} = a^{2} + b^{2}$ for $z = a + bJ$, a non-negative integer, multiplicative because $\Z[J]$ is commutative and the reverse an anti-automorphism; $(J - 1)(J - 1)^{\sim} = 2$. If $J - 1$ had an inverse $z$ in $\Z[J]$ then $2\,z\tilde{z} = 1$ in the integers, which is impossible. If $2m \ge n$ then $2^{m} = J^{m}(J - 1)^{2m}$ is divisible by $(J - 1)^{n}$. Conversely, if $2^{m} = (J - 1)^{n}z$ with $2m < n$, then $J^{m}(J - 1)^{2m} = (J - 1)^{n}z$, and $(J - 1)^{2m}$ is not a zero divisor ($z\tilde{z} = 0$ only for $z = 0$), so $J^{m} = (J - 1)^{n-2m}z$ and $J - 1$ divides the invertible element $J^{m}$, hence is invertible: a contradiction.

\emph{(ii).} The ring is that of clause (6) of \claimref{synthesis.cohomology-of-the-quarter-turn-lens-spaces-from-the-equivariant-cell-structure}. Write $\sigma = \chi - 1$, so it is $A_{n} = \Z[\sigma]/(f, \sigma^{n+1})$ with $f = (1 + \sigma)^{4} - 1 = \sigma(\sigma + 2)(\sigma^{2} + 2\sigma + 2)$, and $\chi^{2} - 1 = \sigma(\sigma + 2)$, $\chi^{2} + 1 = \sigma^{2} + 2\sigma + 2 =: \phi$. In the polynomial ring $\Z[\sigma]$, a unique factorisation domain, $\sigma$ is a prime element and does not divide $\sigma + 2$. Suppose $2^{m}\sigma(\sigma + 2) = Af + B\sigma^{n+1}$ with $A, B \in \Z[\sigma]$. Dividing by $\sigma$: $(\sigma + 2)(2^{m} - A\phi) = B\sigma^{n}$. Since $\sigma$ is prime and does not divide $\sigma + 2$, $\sigma^{n}$ divides $2^{m} - A\phi$: $2^{m} = A\phi + C\sigma^{n}$. Conversely, from $2^{m} = A\phi + C\sigma^{n}$, multiplying by $\sigma(\sigma + 2)$ gives $2^{m}\sigma(\sigma + 2) = Af + C(\sigma + 2)\sigma^{n+1}$. So
\[
2^{m}(1 - \chi^{2}) = 0 \text{ in } A_{n} \iff 2^{m} \in (\phi, \sigma^{n}) \text{ in } \Z[\sigma] \iff 2^{m} = 0 \text{ in } \Z[\sigma]/(\phi, \sigma^{n}).
\]
The polynomial $\phi = (\sigma + 1)^{2} + 1$ is monic, so $\Z[\sigma]/(\phi)$ is the free $\Z$-module on $1, \sigma$ with $(\sigma + 1)^{2} = -1$: it is $\Z[J]$ by $\sigma + 1 = \chi \mapsto J$, and $\sigma \mapsto J - 1$. Hence the condition is that $(J - 1)^{n}$ divides $2^{m}$ in $\Z[J]$, which by (i) is $2m \ge n$. The least such $m$ is $\lceil n/2 \rceil$, and $\lceil n/2 \rceil = \lfloor (n+1)/2 \rfloor$ for integers.

\emph{(iii).} A pointer: the implication is clause (1) of \claimref{synthesis.lens-space-walls-need-more-rungs-than-the-mod-two-height}, and with (ii) and $\xi = \chi^{2}$ it gives $k - 1 \ge \lceil n/2 \rceil$. For even $n$, $\lceil n/2 \rceil = \lfloor n/2 \rfloor$; for odd $n$ it is one more.

\emph{(iv).} Let $N = 2^{a}$, $M = N/2$, $\sigma = \chi - 1$, $w = \chi^{M} - 1 = (1 + \sigma)^{M} - 1$ and $\Phi = \chi^{M} + 1$, so $\chi^{N} - 1 = w\Phi$. The polynomial $w$ has zero constant term and linear coefficient $M \ne 0$, so $w = \sigma w'$ with $w'(0) = M$: the prime element $\sigma$ of $\Z[\sigma]$ does not divide $w'$. As in (ii): $2^{m}w = Aw\Phi + B\sigma^{n+1}$ gives $w'(2^{m} - A\Phi) = B\sigma^{n}$, hence $\sigma^{n}$ divides $2^{m} - A\Phi$; and conversely $2^{m} = A\Phi + C\sigma^{n}$ gives $2^{m}w = Aw\Phi + Cw'\sigma^{n+1}$. So $2^{m}(1 - \chi^{M}) = 0$ in the quotient ring if and only if $(\zeta - 1)^{n}$ divides $2^{m}$ in $R = \Z[\chi]/(\Phi)$, $\zeta$ the class of $\chi$. $R$ is the free $\Z$-module on $1, \zeta, \dots, \zeta^{M-1}$ with $\zeta^{M} = -1$. For $\alpha \in R$ let $d(\alpha)$ be the determinant of multiplication by $\alpha$ on this basis; $d$ is multiplicative, $\alpha$ is invertible in $R$ exactly when $d(\alpha) = \pm 1$ (the adjugate matrix gives an element $\alpha'$ with $\alpha\alpha' = d(\alpha)$), and $\alpha$ is not a zero divisor when $d(\alpha) \ne 0$. Multiplication by $\zeta$ has characteristic polynomial $x^{M} + 1$, so $d(1 - \zeta) = \Phi(1) = 2$ and $d(2) = 2^{M}$. The binomial coefficients $\binom{M}{j}$, $0 < j < M$, are even because $M$ is a power of $2$, so $(x - 1)^{M} = x^{M} + 1 + 2h(x)$ with $h$ an integer polynomial, and in $R$: $(\zeta - 1)^{M} = 2h(\zeta)$. Taking determinants, $(-1)^{M}2^{M} = 2^{M}d(h(\zeta))$, so $d(h(\zeta)) = \pm 1$ and $h(\zeta)$ is invertible: $2 = \varepsilon(\zeta - 1)^{M}$ with $\varepsilon$ invertible. Then, as in (i), $(\zeta - 1)^{n}$ divides $2^{m} = \varepsilon^{m}(\zeta - 1)^{mM}$ exactly when $mM \ge n$: if $mM < n$, cancelling $(\zeta - 1)^{mM}$, which is not a zero divisor, would make $\zeta - 1$ divide an invertible element, while $d(\zeta - 1) = \pm 2$. The least such $m$ is $\lceil n/M \rceil = \lceil 2n/N \rceil$. For $N = 2$ the ring is $\Z$ with $\zeta = -1$; for $N = 4$ it is $\Z[J]$ and this is (ii). The remark on rungs is not part of this result: it would follow from clause (1) of the result named in (iii) applied to a wall whose K-ring is the quotient ring, and that identification for $N \ne 4$ is the classical computation, not stored.
